% !TeX root = ../机器学习笔记.tex
% 由主文档通过 \input 载入；请编译主文档。
\setcounter{chapter}{4} % 保留原书章号。
\chapter{数据处理、模型评估与选择}
\label{chap:ml-evaluation}

评价的对象是从数据得到预测器的整个过程，而不只是最后一个模型对象。数据划分和预处理不当，会使精确实现的算法产生不可靠的结论。

\section{数据来源、采样单位与数据划分}

先明确预测将作用于什么对象。如果目标是新用户，同一用户的多条记录必须作为一个组划分；如果目标是未来时刻，训练应在评价时段之前。逐行随机拆分只有在行与行足够独立且符合目标采样方式时才合适。以比例拆分数据并不自动构成可靠的评价设计。

训练集拟合模型，验证集选择超参数，测试集评价已经固定的流程。\emph{数据泄漏}（data leakage）是训练或选择利用了评价时本不应可用的信息，包括重复实体、未来变量和由全数据统计得到的特征。测试标签尤其不能用于阈值选择、特征筛选或模型取舍；一旦用于这些目的，它便承担了验证功能，需要新的独立评价。

\section{缺失值、编码、尺度与特征工程}

数值特征常作 $z_j=(x_j-\widehat\mu_j)/\widehat s_j$ 标准化，其中均值和标准差只在当前训练数据计算；常数列需要删除或另行规定除零处理。类别特征可用独热编码，类别词表也由训练流程确定。缺失值可用训练均值、中位数或模型估计填补，并视任务保留缺失指示变量；合理性取决于缺失机制，不能认为填补恢复了真实观测。

预处理、特征选择、降维和最终模型应组成同一流程。在交叉验证中，每一折只用该折训练部分拟合这些步骤，再应用于该折验证部分。即使 PCA 不使用标签，预先用全部数据拟合仍可能让验证样本影响表示。目标编码更直接使用标签，应利用训练内的折外编码避免样本标签泄漏到自身特征。

\section{分类与回归的评价指标}

二分类中，真正例、假正例、假反例和真反例分别记为 TP、FP、FN、TN。精确率 $P=\mathrm{TP}/(\mathrm{TP}+\mathrm{FP})$ 衡量预测为正的样本中有多少确实为正；召回率 $R=\mathrm{TP}/(\mathrm{TP}+\mathrm{FN})$ 衡量实际正例中找到了多少。$F_1=2PR/(P+R)$，准确率为 $(\mathrm{TP}+\mathrm{TN})/n$。分母为零时指标未定义，报告中应写清处理约定，多类别的宏平均和微平均也不可混用。

改变概率阈值会改变混淆矩阵。受试者工作特征（receiver operating characteristic，ROC）曲线描绘假正例率与真正例率；精确率--召回率曲线更直接显示正例预测质量与覆盖率的取舍。ROC 曲线下面积可解释为随机正例得分超过随机反例的概率，对并列给半权重；它评价排序，不评价概率校准或某个特定阈值的业务损失。

回归的均方误差为 $\mathrm{MSE}=n^{-1}\sum_i(\widehat y_i-y_i)^2$，平均绝对误差为 $\mathrm{MAE}=n^{-1}\sum_i|\widehat y_i-y_i|$。决定系数
\[
 R^2=1-\frac{\sum_i(y_i-\widehat y_i)^2}{\sum_i(y_i-\bar y)^2}
\]
仅在分母非零时有定义，可能为负，表示所给预测比对评价数据恒预测其均值还差。该评价定义中的 $\bar y$ 不意味着训练时允许获取测试均值。

\section{交叉验证、超参数选择与基线}

$k$ 折交叉验证把开发数据分为 $k$ 份，每次用 $k-1$ 份训练、一份验证，按样本数或既定规则汇总损失。对每个候选超参数重复整个流程，再选择验证表现较好的设置。分类可作分层抽样，相关实体须作分组，时间序列使用滚动或扩展窗口；形式上的折数不能替代任务所需的独立性。

调参后对同一批折报告最好的平均分，仍带有选择偏差。嵌套交叉验证在外层保留评价部分，内层完成所有选择，外层分数评价的是选择程序本身。若保留了独立测试集，也可在开发数据内选择、重新拟合，再进行一次最终测试。基线至少包括训练均值或多数类预测器，以及一个简单模型；所有方法共享数据划分和指标。

\section{不平衡数据、概率校准与评价不确定性}

重采样改变训练样本权重，类别加权改变训练目标，阈值调整改变最终决策，概率校准（probability calibration）修正输出概率与发生率之间的关系；这些操作不能当作同一件事。一个校准良好的二分类概率 $S$ 理想上满足 $P(Y=1\mid S)=S$。校准可利用独立验证集拟合，但用同一数据同时调校准器并报告概率质量会偏乐观。

测试分数有有限样本误差。独立 Bernoulli 正确性指标的样本均值近似标准误为 $\sqrt{\widehat p(1-\widehat p)/n}$，但这不适用于任意相关数据。自助重采样应按用户、序列等真正采样单位进行；验证折因训练集重叠而相关，不能机械当作独立重复。多随机种子只考察算法随机性，不能替代对数据抽样变化的评估。

\section{例题、练习与学习检查}

\begin{example}[准确率与召回率]
一百个测试样本得到 $\mathrm{TP}=8$、$\mathrm{FP}=2$、$\mathrm{FN}=4$、$\mathrm{TN}=86$。评价该分类器。
\end{example}
\begin{solve}
准确率为 $0.94$，精确率为 $0.8$，召回率为 $2/3$，$F_1=8/11$。恒预测负类的准确率已达 $0.88$，故仅看准确率提升会遗漏模型仍错过三分之一正例的事实。是否接受这一取舍还取决于两类错误代价。
\end{solve}

习题：为多用户的时间记录设计一个评价未来新用户的划分；解释先在全数据标准化再交叉验证的问题；为一个有五种阈值候选的分类器安排训练、选阈值与最终测试。实验报告应能让另一人重建划分和选择过程。





